\documentclass{beamer}
\usepackage{../tut-slides}
\usepackage{../mathoperatorsAuD}

\usepackage{lmodern}
\usepackage{amsmath,amssymb}
\usepackage{wasysym}
\usepackage{stmaryrd}
\usepackage{enumerate}
%\usepackage[inline]{enumitem} 		%customize label
%\newcommand{\labelitemi}{\raisebox{1pt}{\scalebox{.9}{$\blacktriangleright$}}}
%\newcommand{\labelitemii}{$\vartriangleright$}
%\newcommand{\labelitemiii}{--}
\setbeamertemplate{itemize item}{\raisebox{1pt}{\scalebox{.9}{$\blacktriangleright$}}}
\setbeamertemplate{itemize subitem}{$\vartriangleright$}

\usepackage{booktabs}
\usepackage{tabularx}
\usepackage{tabu}
\newcommand*\head{\rowfont{\bfseries}}
\newcommand*{\tw}{\rowfont{\ttfamily}}
\renewcommand{\tabularxcolumn}[1]{>{\hspace{0pt}}m{#1}}
\usepackage{multirow}

\usepackage{cancel}

\usepackage{empheq}
\newcommand*\widefbox[1]{\fbox{\hspace{2em} #1 \hspace{2em}}}

\usepackage{tcolorbox}
\newtcolorbox{mymathbox}[1][]{colback=white, sharp corners, #1}

\usepackage{xcolor}
\usepackage{MnSymbol}

\newcommand{\col}[1]{\textcolor{cdpurple}{#1}}
\newcolumntype{R}[1]{>{\centering\arraybackslash}p{#1}}
\usepackage{tabularx}
\renewcommand{\tabularxcolumn}[1]{m{#1}}

\usepackage{ragged2e}
\usepackage{csquotes}

\DeclareMathOperator*{\argmax}{arg\,max}

\DeclareMathOperator{\yield}{yield}
\DeclareMathOperator{\win}{Gewinn}
\DeclareMathOperator{\nowin}{kein Gewinn}
\DeclareMathOperator{\mle}{mle}
\DeclareMathOperator{\rfe}{rfe}

\DeclareMathOperator{\schere}{Schere}
\DeclareMathOperator{\stein}{Stein}
\DeclareMathOperator{\papier}{Papier}



\begin{document}	
	\title{Algorithmen und Datenstrukturen}
	\subtitle{Übung 14: EM-Algorithmus}
	\author{Eric Kunze}
	\email{eric.kunze@mailbox.tu-dresden.de}
	\city{TU Dresden}
%	\institute{Lehrstuhl für Grundlagen der Programmierung}
	\titlegraphic{\includegraphics[width=2cm]{../TUD-white.pdf}}
	\date{24.01.2022}

	\maketitle

%%%%%%%%%%%%%%%%%%%%%%%%%%%%%%%%%%%%%%%%%%%%%%%%%%%%%%%%%%%%%%%%%%%%%%%%%%%%%%%%%%%

\section{Aufgabe 1}

%\begin{frame} \frametitle{Wahrscheinlichkeitstheorie}
%	\justifying \footnotesize
%	\setbeamerfont{itemize/enumerate body}{size={\footnotesize}}
%	\setbeamerfont{itemize/enumerate subbody}{size={\footnotesize}}
%	
%	Wir betrachten ein Zufallsexperiment mit Ergebnismenge $X$ und einer Wahrscheinlichkeitsverteilung $p$. Eine Menge $\mathcal{M}$ von Wahrscheinlichkeitsverteilungen heißt Wahrscheinlichkeitsmodell.
%	
%	\textbf{Praktisches Problem:} \\
%	nur \textit{beobachtete} Ergebnisse verfügbar; genaue Verteilung unbekannt
%	
%	\begin{itemize}
%		\item \textbf{$X$-Korpus} $h : X \to \R^{\ge 0}$ 
%		\begin{itemize}
%			\item $h$ zählt beobachtete Ergebnisse; absolute Häufigkeiten
%		\end{itemize}
%		\item Wie gut passt eine Verteilung $p$ zu einem Korpus $h$?
%		\begin{itemize}
%			\item Likelihood von $h$ unter $p$:
%			\begin{equation*}
%				L(h,p) = \prod_{x \in X} p(x)^{h(x)} .
%			\end{equation*}
%		\end{itemize}
%		\item \textbf{Schätzung der Verteilung} aus beobachteten Ergebnissen
%		\begin{itemize}
%			\item wähle die am besten passendste Verteilung
%			\item Maximum-Likelihood-Schätzer: 
%			\begin{equation*}
%				\operatorname{mle}(h,\mathcal{M}) = \argmax_{p \in \mathcal{M}} L(h,p) .
%			\end{equation*}
%			\item Falls $\mathcal{M} = \mathcal{M}(X)$: $\mle(h, \mathcal{M}) = \rfe(h)$
%		\end{itemize}
%	\end{itemize}
%\end{frame}

\begin{frame} \frametitle{Wahrscheinlichkeitstheorie}
	\justifying \footnotesize
	Wir betrachten ein Zufallsexperiment $(X,p)$ mit
	\begin{itemize}
		\item Ergebnismenge $X$ und
		\item einer Funktion $p \colon X \to [0,1]$ mit $\sum_{x \in X} p(x) = 1$ (\textbf{Wahrscheinlichkeitsverteilung von $X$})
	\end{itemize}
	\pause
	
	Die Menge aller Wahrscheinlichkeitsverteilungen über $X$ sei $\mathcal{M}(X)$. Jede Teilmenge $\mathcal{M} \subseteq \mathcal{M}(X)$ heißt \textbf{Wahrscheinlichkeitsmodell}.
	\pause
	
	Ein Wahrscheinlichkeitsmodell $\mathcal{M}$ heißt \textbf{beschränkt}, falls $\mathcal{M} \neq \mathcal{M}(X)$; andernfalls unbeschränkt.
	\pause
	
	Führen wir nun zwei Zufallsexperimente nacheinander aus und nehmen dabei an, dass die beiden Experimente \textit{unabhängig} voneinander sind. Folge das erste Experiment einer Verteilung $p_1 \in \mathcal{M}(X_1)$ und das zweite Experiment einer Verteilung $p_2 \in \mathcal{M}(X_2)$, dann ist $p_1 \times p_2 \in \mathcal{M}(X_1 \times X_2)$ eine Verteilung auf der Ergebnismenge $X_1 \times X_2$ unseres zweistufigen Experiments:
	\vspace{-.5\baselineskip}
	\begin{equation*}
		(p_1 \times p_2)(a,b) = p_1(a) * p_2(b) .
	\end{equation*}
	\vspace{-1.5\baselineskip}
	\begin{center}
		\enquote{Einzelwahrscheinlichkeiten multiplizieren / erste Pfadregel}
	\end{center}
\end{frame}

\begin{frame} \frametitle{Korpora und Korpuswahrscheinlichkeiten}
	\justifying \footnotesize
	Oftmals wissen wir aber die zugrundeliegende Verteilung nicht, sondern können lediglich die Ergebnisse des Experiments wahrnehmen. Zählen wir diese Beobachtungen, dann nennen wir das einen \textbf{$X$-Korpus} modelliert durch eine Funktion $h \colon X \to \R^{\ge 0}$.
	Man definiert die \textbf{Korpuswahrscheinlichkeit} / \textbf{Likelihood} von $h$ unter einer Verteilung $p$ als
	\begin{equation*}
		L(h,p) = \prod_{x \in X} p(x)^{h(x)} .
	\end{equation*} 
	Nun kennen wir aber die Verteilung $p$ nicht und müssen sie daher aus den beobachteten Daten schätzen. Dies macht der\textbf{ Maximum-Likelihood-Schätzer} (MLE)
	\begin{equation*}
		\operatorname{mle}(h,\mathcal{M}) = \argmax_{p \in \mathcal{M}} L(h,p) .
	\end{equation*}
	Solange das Modell unbeschränkt gewählt wird, d.h. es werden alle Verteilungen über $X$ zugelassen, dann wird der MLE zur relativen Häufigkeit von $h$.
\end{frame}

\begin{frame} \frametitle{Unvollständige Daten}
	\justifying \footnotesize
	Bisher sind wir davon ausgegangen, dass die Daten stets vollständig waren, d.h. wir konnten jedes Ergebnis beobachten. In der Realität können aber oftmals nur Gruppen von Ergebnissen beobachtet werden; z.B. gewinne oder verliere ich bei einem Spiel. Wir wissen aber nicht, welches Ergebnis genau erzielt wurde. 
	
	Sei $Y$ die Menge der Beobachtungen. Die Beobachtungsfunktion $\operatorname{yield} \colon X \to Y$ ordnet jedem Ergebnis seine Beobachtung zu. Die Umkehrabbildung ordnet dann jeder Beobachtung eine Menge von möglichen Ergebnissen zu, die zu dieser Beobachtung führen, d.h.
	\begin{equation*}
		A \colon Y \to \mathcal{P}(X) \quad \text{mit} \quad A(y) = \menge{x \in X: \operatorname{yield}(x) = y} .
	\end{equation*}
	Diese Funktion heißt \textbf{Analysator}. 
	\pause
	
	Sei $h$ ein $Y$-Korpus, d.h. $h$ zählt Beobachtungen (nicht Ergebnisse). Die \textbf{Korpuswahrscheinlichkeit} / \textbf{Likelihood} von $h$ unter einer Verteilung $p$ ist 
	\begin{equation*}
		L(h,p) = \prod_{y \in Y} \left( \sum_{x \in A(y)} p(x) \right)^{h(y)}.
	\end{equation*}
	Der MLE bleibt wie er war: $\operatorname{mle}(h,\mathcal{M}) = \argmax_{p \in \mathcal{M}} L(h,p)$.
\end{frame}

\begin{frame} \frametitle{Aufgabe 1}
	\small
	Bestimmen Sie für die folgenden Szenarien die Menge $X$ der Ergebnisse und die Menge $Y$ der
	Beobachtungen. Bestimmen Sie außerdem den Analysator.
	\begin{enumerate}[(a)]
		\item Werfen zweier unabhängiger Münzen. Sie können nur beobachten, ob beide Münzen dieselbe oder verschiedene Seiten zeigen.
		\item Werfen zweier Würfel, wobei Sie nur die Summe der Augenzahlen beobachten.
		\item Zwei Spieler spielen Schere-Stein-Papier. Sie beobachten lediglich, welcher Spieler gewonnen hat bzw. ob das Spiel unentschieden ausging.
	\end{enumerate}
\end{frame}

\begin{frame} \frametitle{Aufgabe 1}
	\footnotesize
	\begin{enumerate}[(a)]
		\item \textbf{zweimaliger Münzwurf -- Beobachtung der Gleichheit}
		\begin{equation*}
			X = \menge{K,Z} \times \menge{K,Z} \qquad \und \qquad
			Y = \menge{\operatorname{gleich}, \operatorname{ungleich}} 
		\end{equation*}
		\pause
		Der Analysator ordnet jeder Beobachtung $y \in Y$ die Menge der Ergebnisse aus $X$ zu, die zur Beobachtung $y$ führen, also
		\begin{align*}
			A(\operatorname{gleich}) &= \menge{(K,K), (Z,Z)} \\
			A(\operatorname{ungleich}) &= \menge{(K,Z), (Z,K)}.			
		\end{align*}
		\pause
		\item \textbf{zweimaliger Würfelwurf -- Beobachtung der Augensumme}
		\begin{equation*}
			X = \menge{1, \dots, 6} \times \menge{1, \dots, 6}
			\qquad \und \qquad
			Y = \menge{2, \dots, 12} 
		\end{equation*}
		\pause
		Analysator: $A(x) = \menge{(i,j) \in X : i + j = x}$, d.h. konkret
		\begin{align*}
			A(2) &= \menge{(1,1)} \\
			A(3) &= \menge{(1,2), (2,1)} \\
			A(4) &= \menge{(1,3), (3,1), (2,2)} \\
			&\vdots \\
			A(12) &= \menge{(6,6)}			
		\end{align*}
	\end{enumerate}
\end{frame}

\begin{frame} \frametitle{Aufgabe 1}
	\footnotesize
	\begin{enumerate}[(a)]
		\setcounter{enumi}{2}
		\item \textbf{Schere, Stein, Papier -- Beobachtung des Gewinners}
		\begin{align*}
			X &= \menge{\operatorname{Schere}, \operatorname{Stein}, \operatorname{Papier}}^2 \\
			Y &= \menge{\operatorname{Spieler1}, \operatorname{Spieler2}, \operatorname{Unentschieden}} 
		\end{align*}
		\pause
		
		Analysator:
		\begin{align*}
			A(\operatorname{Spieler1}) &= \menge{(\schere, \papier), (\stein, \schere), (\papier, \stein)} \\
			A(\operatorname{Spieler2}) &= \menge{(\papier, \schere), (\schere, \stein), (\stein, \papier)} \\
			A(\operatorname{Unentschieden}) &= \menge{(\papier, \papier), (\stein, \stein), (\schere, \schere)}		
		\end{align*}
	\end{enumerate}
\end{frame}

\section{Aufgabe 2}

\begin{frame} \frametitle{Problem: Ergebnisse vs. Beobachtungen}
	\justifying \footnotesize 
	\makebox[2cm][l]{\textbf{Gegeben:}} ein $Y$-Korpus von \textit{Beobachtungen} \\	
	\makebox[2cm][l]{\textbf{Gesucht:}} eine Verteilung $p \in \mathcal{M} \subseteq \mathcal{M}(X)$ auf \textit{Ergebnissen}
	
	\bigskip \pause
	
	\textbf{Problem:} Wir wollen die Verteilung der \textit{Ergebnisse} schätzen. Um den MLE nutzen zu können, brauchen wir dafür einen $X$-Korpus. Jedoch ist uns nur ein $Y$-Korpus gegeben, da wir nur \textit{Beobachtungen} wahrnehmen können. 
	\pause
	
	\textbf{Ausweg:} Erweiterung des $Y$-Korpus $h$ zu einem $X$-Korpus $h_1$ auf vollständigen Daten
	\begin{equation*}
		h_i(x) = h(\operatorname{yield}(x)) \cdot \frac{p_{i-1}(x)}{\sum_{x' \in A(\operatorname{yield}(x))} p_{i-1}(x)} \qquad \text{ für alle } x \in X
	\end{equation*}
	Dazu benötigen wir eine gewisse Vorkenntnis mit der Verteilung $p_{i-1}$, die wir aus dem vorherigen Iterationsschritt bzw. einer initialen Vermutung bekommen.
	
	Dies ist der E-Schritt des EM-Algorithmus.
\end{frame}

\begin{frame} \frametitle{EM-Algorithmus}
	\footnotesize
	\textbf{Input:} Analysator $A : Y \to \mathcal{P}(X)$, Modell $\mathcal{M} \subseteq \mathcal{M}(X)$, \\
	\begin{itemize}
		\item $Y$-Korpus $h$ von Beobachtungen,
		\item initiale Wahrscheinlichkeitsverteilung $p_0$ 
	\end{itemize}
	
	Eine Iteration des Algorithmus besteht aus den folgenden beiden Schritten:
	\begin{description}
		\item[E-Schritt] \textit{Expectation} \\
		Bestimmte die versteckten Eigenschaften mithilfe der Parameter aus der vorherigen Iteration.
		\begin{equation*}
		h_i(x) = h(\operatorname{yield}(x)) \cdot \frac{p_{i-1}(x)}{\sum_{x' \in A(\operatorname{yield}(x))} p_{i-1}(x)}
		\end{equation*}
		\item[M-Schritt] \textit{Maximization} \\
		Bestimmte die neuen Parameter mithilfe der vollständigen Eigenschaften aus dem E-Schritt.
		\begin{equation*}
		p_i = \argmax_{p \in \mathcal{M}} L(h_i, p)
		\end{equation*}
	\end{description}
\end{frame}

\begin{frame}\frametitle{Ein weiteres Problem?}
	\justifying \footnotesize 
	\textbf{Problem:} Im M-Schritt bleibt ein MLE zu berechnen, der für beschränkte Modelle $\mathcal{M} \neq \mathcal{M}(X)$ nicht leicht zu bekommen ist. Eine Vereinfachung ist aber für mehrstufige Zufallsversuche möglich, die als \textit{unabhängig} vorausgesetzt werden.
	\pause
	
	Haben wir zwei unabhängige Teilversuche mit Ergebnismengen $X_1$ und $X_2$, dann soll die Verteilung die Unabhängigkeit widerspiegeln und wir verwenden daher das Modell 
	\begin{equation*}
		\mathcal{M} = \menge{p^1 \times p^2 : p^1 \in \mathcal{M}(X_1), p^2 \in \mathcal{M}(X_2)} \neq \mathcal{M}(X_1 \times X_2),
	\end{equation*}
	der MLE wird also nicht zur relativen Häufigkeit auf $X_1 \times X_2$. 
	\pause
	
	\textbf{Ausweg:} Die Unabhängigkeit rettet uns: wir erhalten den MLE bzgl. $\mathcal{M}$ als unabhängiges Produkt der relativen Häufigkeiten \textit{auf den Teilexperimenten}. Dazu ist ein ($X_1 \times X_2$)-Korpus $h$ nicht unbedingt hilfreich. Wir brauchen vielmehr die Teilkorpora $h^1$ auf $X_1$ und $h^2$ auf $X_2$. Dann gilt
	\begin{equation*}
		\mle(h, \mathcal{M}) = \rfe(h^1) \times \rfe(h^2) \textcolor{cdgray!70}{ \neq \rfe(h)}.
	\end{equation*}
	\pause
	
	\textbf{Problem:} Wie bekommen wir $h^1$ und $h^2$? \pause  \ --- \
	\textbf{Ausweg:} Marginalisierung
\end{frame}

\begin{frame} \frametitle{Marginalisierung}
	\justifying \footnotesize
	Wir betrachten die zwei Ergebnismengen $X_1$ und $X_2$. Das Modell sei gegeben durch das unabhängige Produkt der Modelle auf $X_1$ und der Modelle auf $X_2$, d.h. $\mathcal{M} = \menge{ p^1 \times p^2 : p^1 \in \mathcal{M}(X_1), p^2 \in \mathcal{M}(X_2)}$. Weiter sei $h$ ein $X_1 \times X_2$-Korpus. Die Teilkorpora $h^1$ auf $X_1$ und $h^2$ auf $X_2$ erhalten wir durch \textbf{Marginalisierung}
	\begin{equation*}
		\begin{aligned}
			h^1(x_1) &= \sum_{x_2 \in X_2} h(x_1, x_2) & \text{für alle } x_1 &\in X_1 \\
			h^2(x_2) &= \sum_{x_1 \in X_1} h(x_1, x_2) & \text{für alle } x_2 &\in X_2 \\
		\end{aligned}
	\end{equation*}
	\pause
	Die Summen entsprechen dabei gerade Zeilen- bzw. Spaltensummen, wenn man $h$ in einer Tabelle notiert.
	\begin{equation*}
		\begin{array}{|c||ccc|c|}
		\hline
			X_1 \backslash X_2 & \alpha & \dots & \omega & \\ \hline \hline
			a & h(a,\alpha) & \dots & h(a,\omega) & h^1(a) \\
			\vdots &\vdots & \ddots & \vdots & \vdots \\
			z & h(z,\alpha) & \dots & h(z,\omega) & h^1(z) \\ \hline
			& h^2(\alpha) & \dots & h^2(\omega) & \\ \hline
		\end{array}
	\end{equation*}	
\end{frame}

\begin{frame} \frametitle{Aufgabe 2 --- Teil (a)}
	\justifying \small	
	Das Spiel wird gewonnen, wenn beide Münzen auf der gleichen Seite landen.
	
	\pause
	
	Damit ist der Analysator $A \colon \menge{\win, \nowin} \to \mathcal{P}(X)$ gegeben durch
	\begin{equation*}
		\begin{aligned}
			A(\win) &= \menge{x \in X: \yield(x) = \win} \\
					&= \menge{(K,K), (Z,Z)} \\
			A(\nowin) &= \menge{x \in X: \yield(x) = \nowin} \\
					  &= \menge{(K,Z), (Z,K),(R,K),(R,Z)}
		\end{aligned}
	\end{equation*}
\end{frame}

\begin{frame} \frametitle{Aufgabe 2 --- Teil (b)}
	\justifying \small 
	Wir können nur die Beobachtungen $\win$ und $\nowin$ feststellen. 
	
	Wir spielen das Spiel $24$ Mal und gewinnen $6$ Mal. Gesucht ist nun der $Y$-Korpus $h$, d.h. wie oft beobachten wir die Ereignisse $\win$ und $\nowin$.
	\pause
	\begin{equation*}
		h(\win) = 6 \qquad\qquad h(\nowin) = 18
	\end{equation*}
\end{frame}

\begin{frame} \frametitle{Aufgabe 2 --- Teil (c)}
	\justifying \footnotesize
	Gegeben ist nun eine initiale Wahrscheinlichkeitsverteilung $q_0 = q_0^1 \times q_0^2$ über den vollständigen Daten mit
	\begin{align*}
	q_0^1(K) &= \frac{2}{5} & q_0^2(K) &= \frac{1}{3} \\
	q_0^1(R) &= \frac{1}{5} \\ \pause
	\Rightarrow
	q_0^1(Z) &= 1 - q_0^1(K) - q_0^1(R) = \frac{2}{5} 
	&
	q_0^2(Z) &= 1 - q_0^1(K) = \frac{2}{3}
	\end{align*}
	
	\onslide<2->{
	Mit dem unabhängigen Produkt erhalten wir
	\begin{align*}
	q_0(K,K) &= q_0^1(K) \cdot q_0^2(K) = \frac{2}{15} 
	& q_0(K,Z) &= q_0^1(K) \cdot q_0^2(Z) = \frac{4}{15} \\
	q_0(Z,K) &= q_0^1(Z) \cdot q_0^2(K) = \frac{2}{15}
	& q_0(Z,Z) &= q_0^1(Z) \cdot q_0^2(Z) = \frac{4}{15} \\
	q_0(R,K) &= q_0^1(R) \cdot q_0^2(K) = \frac{1}{15}
	& q_0(R,Z) &= q_0^1(R) \cdot q_0^2(Z) = \frac{2}{15}
	\end{align*}
	}
\end{frame}

\begin{frame} \frametitle{Aufgabe 2 --- Teil (c)}
	\justifying \footnotesize
	\textbf{E-Schritt:} Erweiterung von $h$ auf $h_1$ mit folgender Formel:
	\begin{equation*}
		h_1(x) = h(\yield(x)) \cdot \frac{q_0(x)}{\sum\limits_{x' \in A(\yield(x))} q_0(x')}
	\end{equation*}
	\pause
	Damit ergibt sich dann zum Beispiel für das Ergebnis $(K,K)$
	\begin{align*}
		h_1(K,K) &= h(\win) \cdot \frac{q_0(K,K)}{\sum\limits_{x' \in \menge{(K,K),(Z,Z)}} q_0(x')} \\
		&= h(\win) \cdot \frac{q_0(K,K)}{q_0(K,K) + q_0(Z,Z)} \\
		&= 6 \cdot \frac{\frac{2}{15}}{\frac{2}{15} + \frac{4}{15}} \\
		&= 2
	\end{align*}
	\pause
	Mit gleicher Rechnung erhalten wir für die restlichen Ereignisse
	\begin{align*}
		 && h_1(Z,K) &= 4 & h_1(R,K) &= 2 \\
		h_1(K,Z) &= 8 & h_1(Z,Z) &= 4 & h_1(R,Z) &= 4
	\end{align*}
\end{frame}

\begin{frame} \frametitle{Aufgabe 2 --- Teil (d)}
	\justifying \footnotesize
	\textbf{M-Schritt:} Bestimmung der Teilkorpora $h_1^1$ bzw. $h_1^2$ durch \textit{Marginalisierung}:
	\pause
	\begin{equation*}
		\begin{array}{|c||cc|c|}
		\hline
		X_1 \backslash X_2 & K & Z & \\ \hline \hline
		K & h_1(K,K) & h_1(K,Z) & h_1^1(K) \\
		Z & h_1(Z,K) & h_1(Z,Z) & h_1^1(Z) \\
		R & h_1(R,K) & h_1(R,Z) & h_1^1(R) \\ \hline
		& h_1^2(K) & h_1^2(Z) & \\ \hline
		\end{array}
		\qquad \leadsto \qquad 
		\begin{array}{|c||cc|c|}
		\hline
		X_1 \backslash X_2 & K & Z & \\ \hline \hline
		K & 2 &  8 & 10 \\
		Z & 4 &  4 &  8 \\
		R & 2 &  4 &  6 \\ \hline
		& 8 & 16 & 24 \\ \hline
		\end{array}
	\end{equation*}
\end{frame}

\begin{frame} \frametitle{Aufgabe 2 --- Teil (e)}
	\justifying \footnotesize
	Nun bestimmen wir noch die relativen Häufigkeiten mit der Formel
	\begin{equation*}
		\rfe(h)(x) \defeq \frac{h(x)}{\abs{h}} \quad \mit \quad \abs{h} \defeq \sum_{x \in X} h(x)
	\end{equation*}
	\pause
	Wenden wir dies nun auf $h_1$ und $h_2$ an, so erhalten wir
	\begin{align*}
		q_1^1(K) = \rfe(h_1^1)(K) &= \frac{h_1^1(K)}{h_1^1(K) + h_1^1(Z) + h_1^1(R)} = \frac{10}{24} = \frac{5}{12} \\
		q_1^1(Z) = \rfe(h_1^1)(Z) &= \frac{h_1^1(Z)}{h_1(K) + h_1^1(Z) + h_1^1(R)} = \frac{8}{24} = \frac{1}{3} \\
		q_1^1(R) = \rfe(h_1^1)(R) &= \frac{h_1^1(R)}{h_1^1(K) + h_1^1(Z) + h_1^1(R)} = \frac{6}{24} = \frac{1}{4} \\
		\intertext{und}
		q_1^2(K) = \rfe(h_1^2)(K) &= \frac{h_1^2(K)}{h_1^2(K) + h_1^2(Z)} = \frac{8}{24} = \frac{1}{3} \\
		q_1^2(Z) = \rfe(h_1^2)(Z) &= \frac{h_1^2(Z)}{h_1^2(K) + h_1^2(Z)} = \frac{16}{24} = \frac{2}{3} \\
	\end{align*}
\end{frame}

\end{document}
